\section{Cálculos Estatísticos}

A FX-82MS dispõe de um modo dedicado a cálculos estatísticos com uma
variável, denominado Modo SD (\textit{Standard Deviation}). Nesse modo, a
calculadora mantém um editor interno de dados e disponibiliza funções para
o cálculo automático de média, desvio padrão e outras grandezas derivadas,
eliminando a necessidade de realizar essas operações manualmente etapa por
etapa \cite{casio2004}.

Para ativar o Modo SD, pressione \textbf{MODE} seguido de \textbf{2}. O
visor exibirá o indicador \textbf{SD} no topo, confirmando que a calculadora
está pronta para receber dados estatísticos \cite{casio2004}. Recomenda-se
limpar os dados de sessões anteriores antes de iniciar uma nova entrada,
utilizando a sequência \textbf{SHIFT} $\to$ \textbf{CLR} $\to$ \textbf{1}
(Scl), que apaga todos os dados do editor sem alterar as demais configurações
da calculadora.

\subsection{Inserindo e Editando Dados}

No Modo SD, os dados são inseridos um a um por meio da tecla \textbf{DT},
que corresponde à função secundária da tecla \textbf{M+}. Cada valor deve
ser digitado e confirmado com \textbf{DT}; o visor exibe brevemente o número
de elementos inseridos até o momento, identificado no visor como \textbf{n}
\cite{casio2004}.

Para inserir os valores 4, 7 e 9, por exemplo, a sequência é:

\begin{table}[H]
    \centering
    \begin{tabular}{lll}
        \hline
        \textbf{Etapa} & \textbf{Entrada} & \textbf{Visor} \\ \hline
        Ativar Modo SD & \texttt{MODE 2} & SD \\
        Limpar dados & \texttt{SHIFT CLR 1} & 0 \\
        Inserir primeiro valor & \texttt{4 DT} & n = 1 \\
        Inserir segundo valor & \texttt{7 DT} & n = 2 \\
        Inserir terceiro valor & \texttt{9 DT} & n = 3 \\ \hline
    \end{tabular}
\end{table}

Caso um valor tenha sido inserido incorretamente, é possível removê-lo
sem apagar todo o conjunto. Para excluir uma entrada, redigite o valor
incorreto e pressione \textbf{SHIFT} seguido de \textbf{DT} (função
\textbf{CD} -- \textit{Count Down}), que remove a última ocorrência daquele
valor do editor \cite{casio2004}. Em seguida, insira o valor correto
normalmente com \textbf{DT}.

É importante observar que a FX-82MS não permite visualizar ou navegar
pela lista de dados inseridos. Por isso, recomenda-se anotar os valores em
papel à medida que são inseridos, facilitando a identificação e correção de
eventuais erros.

\subsection{Média e Soma}

Após a inserção dos dados no Modo SD, a FX-82MS disponibiliza as
grandezas estatísticas por meio da tecla \textbf{SHIFT} seguida de
\textbf{2} (função \textbf{S-VAR}), que abre um menu com as variáveis
estatísticas disponíveis \cite{casio2004}. A média aritmética e as somas
acumuladas são acessadas por esse menu e exibidas no visor ao confirmar
com \textbf{=}.

As grandezas relacionadas à média e à soma são \cite{casio2004}:

\begin{itemize}
    \item $\bar{x}$ \textbf{(S-VAR 1)} -- média aritmética dos valores inseridos,
    calculada automaticamente pela calculadora como:
    \begin{equation}
        \bar{x} = \frac{\sum x_i}{n}
    \end{equation}
    \item $\sum x$ \textbf{(S-VAR 4)} -- soma de todos os valores inseridos;
    \item $\sum x^2$ \textbf{(S-VAR 5)} -- soma dos quadrados de todos os valores inseridos;
    \item $n$ \textbf{(S-VAR 6)} -- número total de elementos inseridos.
\end{itemize}

Para recuperar qualquer uma dessas grandezas, pressione \textbf{SHIFT}
$\to$ \textbf{2} (S-VAR), selecione o número correspondente à grandeza
desejada e confirme com \textbf{=}. Considerando o conjunto \{4, 7, 9\}
já inserido:

\begin{table}[H]
    \centering
    \begin{tabular}{lll}
        \hline
        \textbf{Grandeza} & \textbf{Entrada} & \textbf{Resultado} \\ \hline
        Média $\bar{x}$ & \texttt{SHIFT 2 1 =} & 6{,}666\ldots \\
        Soma $\sum x$ & \texttt{SHIFT 2 4 =} & 20 \\
        Soma $\sum x^2$ & \texttt{SHIFT 2 5 =} & 146 \\
        Total $n$ & \texttt{SHIFT 2 6 =} & 3 \\ \hline
    \end{tabular}
\end{table}

O valor de $\bar{x}$ retornado pela calculadora é mantido internamente
com precisão completa, mesmo que o visor exiba uma versão arredondada.
Ao utilizar o resultado de $\bar{x}$ em cálculos subsequentes por meio
da tecla \textbf{Ans}, a precisão total é preservada, evitando erros de
arredondamento acumulado \cite{casio2004}.

\subsection{Desvio Padrão}

O desvio padrão é uma medida de dispersão que indica o quanto os valores
de um conjunto se afastam, em média, da média aritmética. A FX-82MS calcula
automaticamente duas versões do desvio padrão no Modo SD, acessíveis pelo
mesmo menu \textbf{S-VAR} utilizado para a média e as somas
\cite{casio2004}.

A distinção entre as duas versões está na natureza do conjunto de dados
analisado: se os dados representam a população completa, utiliza-se o desvio
padrão populacional; se representam uma amostra extraída de uma população
maior, utiliza-se o desvio padrão amostral \cite{casio2004}.

\subsubsection{Desvio Padrão Populacional}

O desvio padrão populacional $s x$ considera que os dados inseridos
representam todos os elementos da população. Sua fórmula é \cite{casio2004}:

\begin{equation}
    s_1 = \sqrt{\frac{\sum (x_i - \bar{x})^2}{n}}
\end{equation}

É acessado por meio de \textbf{SHIFT} $\to$ \textbf{2} (S-VAR) $\to$
\textbf{2}. Considerando o conjunto \{4, 7, 9\}:

\begin{table}[H]
    \centering
    \begin{tabular}{lll}
        \hline
        \textbf{Grandeza} & \textbf{Entrada} & \textbf{Resultado} \\ \hline
        Desvio padrão populacional $s_1$ & \texttt{SHIFT 2 2 =} & $\approx 2{,}055$ \\ \hline
    \end{tabular}
\end{table}

\subsubsection{Desvio Padrão Amostral}

O desvio padrão amostral $s_2$ considera que os dados inseridos representam
uma amostra da população, aplicando a correção de Bessel no denominador para
produzir uma estimativa não tendenciosa \cite{casio2004}:

\begin{equation}
    s_2 = \sqrt{\frac{\sum (x_i - \bar{x})^2}{n - 1}}
\end{equation}

É acessado por meio de \textbf{SHIFT} $\to$ \textbf{2} (S-VAR) $\to$
\textbf{3}. Para o mesmo conjunto \{4, 7, 9\}:

\begin{table}[H]
    \centering
    \begin{tabular}{lll}
        \hline
        \textbf{Grandeza} & \textbf{Entrada} & \textbf{Resultado} \\ \hline
        Desvio padrão amostral $s_2$ & \texttt{SHIFT 2 3 =} & $\approx 2{,}517$ \\ \hline
    \end{tabular}
\end{table}

A diferença entre $s_1$ e $s_2$ é mais pronunciada em conjuntos
pequenos e tende a diminuir à medida que $n$ cresce. Em contextos acadêmicos,
o enunciado do problema geralmente indica qual versão utilizar; na ausência
dessa informação, recomenda-se verificar se os dados representam uma
população completa ou uma amostra antes de escolher a função adequada.

\subsection{Outras Grandezas Disponíveis}

Além da média e dos desvios padrão apresentados nas seções anteriores, o
menu S-VAR reúne todas as grandezas calculadas automaticamente pela
calculadora a partir dos dados inseridos no Modo SD. A tabela a seguir
apresenta uma visão consolidada das seis grandezas disponíveis, com os
respectivos atalhos e significados \cite{casio2004}.

\begin{table}[H]
    \centering
    \begin{tabular}{lll}
        \hline
        \textbf{Atalho} & \textbf{Grandeza} & \textbf{Significado} \\ \hline
        \texttt{SHIFT 2 1} & $\bar{x}$    & Média aritmética \\
        \texttt{SHIFT 2 2} & $s_1$        & Desvio padrão populacional \\
        \texttt{SHIFT 2 3} & $s_2$        & Desvio padrão amostral \\
        \texttt{SHIFT 2 4} & $\sum x$     & Soma dos valores \\
        \texttt{SHIFT 2 5} & $\sum x^2$   & Soma dos quadrados dos valores \\
        \texttt{SHIFT 2 6} & $n$          & Número de elementos \\ \hline
    \end{tabular}
\end{table}

Os resultados do menu S-VAR podem ser utilizados diretamente em cálculos
subsequentes por meio da tecla \textbf{Ans} \cite{casio2004}, permitindo
derivar grandezas que não constam no menu, como a variância e o coeficiente
de variação. A variância corresponde ao quadrado do desvio padrão e pode ser
obtida elevando o resultado retornado pelo menu ao quadrado com a tecla $x^2$:

\begin{table}[H]
    \centering
    \begin{tabular}{lll}
        \hline
        \textbf{Etapa} & \textbf{Entrada} & \textbf{Resultado} \\ \hline
        Recuperar $s_2$ & \texttt{SHIFT 2 3 =}  & valor de $s_2$ \\
        Variância amostral $s_2^2$ & \texttt{Ans x\textsuperscript{2} =} & $s_2^2$ \\ \hline
    \end{tabular}
\end{table}

Para o conjunto \{4, 7, 9\} utilizado nas seções anteriores, o valor
$s_2 \approx 2{,}517$ resulta em variância amostral $s_2^2 \approx 6{,}333$.
De modo análogo, a variância populacional obtida a partir de $s_1 \approx 2{,}055$
é $s_1^2 \approx 4{,}222$.

\subsection{Dados com Frequência}

Quando os dados provêm de uma tabela de frequências ou contêm valores
repetidos, a FX-82MS permite inserir cada valor acompanhado de sua frequência
absoluta, evitando a digitação individual de cada ocorrência. Esse recurso é
especialmente útil em exercícios de Estatística que apresentam os dados já
resumidos em tabelas \cite{casio2004}.

No Modo SD, a entrada de um valor $x$ com frequência $f$ é realizada pela
sequência \cite{casio2004}:

\begin{center}
    \textit{valor} \textbf{;} \textit{frequência} \textbf{DT}
\end{center}

Onde o separador \textbf{;} é acessado pela tecla correspondente do teclado,
conforme indicado no painel da calculadora. Ao confirmar com \textbf{DT}, o
visor exibe o número total acumulado de elementos $n$, já computando a
frequência informada. Quando a frequência é igual a 1, o separador pode ser
omitido e o valor inserido diretamente com \textbf{DT}.

Considere o conjunto de tempos de atendimento \{2, 4, 4, 4, 5, 5, 7, 9\}, cujos valores repetidos tornam a
entrada por frequência vantajosa. Sua tabela de frequências é:

\begin{table}[H]
    \centering
    \begin{tabular}{ll}
        \hline
        \textbf{Valor} ($x_i$) & \textbf{Frequência} ($f_i$) \\ \hline
        2 & 1 \\
        4 & 3 \\
        5 & 2 \\
        7 & 1 \\
        9 & 1 \\ \hline
    \end{tabular}
\end{table}

A inserção com frequência reduz o número de entradas de 8 para 5:

\begin{table}[htbp]
    \centering
    \begin{tabular}{lll}
        \hline
        \textbf{Etapa} & \textbf{Entrada} & \textbf{Visor} \\ \hline
        Ativar Modo SD      & \texttt{MODE 2}       & SD    \\
        Limpar dados        & \texttt{SHIFT CLR 1}  & 0     \\
        Valor 2, freq.\ 1  & \texttt{2 DT}         & n = 1 \\
        Valor 4, freq.\ 3  & \texttt{4 ; 3 DT}     & n = 4 \\
        Valor 5, freq.\ 2  & \texttt{5 ; 2 DT}     & n = 6 \\
        Valor 7, freq.\ 1  & \texttt{7 DT}         & n = 7 \\
        Valor 9, freq.\ 1  & \texttt{9 DT}         & n = 8 \\ \hline
    \end{tabular}
\end{table}

Após a inserção, as grandezas estatísticas são recuperadas normalmente pelo
menu S-VAR, com $n = 8$ elementos computados de acordo com as frequências
informadas. O conjunto inserido por frequência é equivalente, para fins de
cálculo, à inserção individual de todos os elementos \cite{casio2004}.

Para remover uma entrada inserida com frequência incorreta, redigite o valor
e a frequência no mesmo formato e pressione \textbf{SHIFT} seguido de
\textbf{DT} (função \textbf{CD}).

\subsection{Exemplos Práticos}

Esta subseção apresenta aplicações de cálculos estatísticos com
exemplos contextualizados em Probabilidade e Estatística.

\subsubsection*{Exemplo 1 — Distribuição Binomial}

Em um processo de fabricação, a probabilidade de uma peça apresentar defeito é $p = 0{,}2$ ($20\%$). Se $5$ peças forem selecionadas aleatoriamente, calcule a probabilidade de exatamente $2$ peças serem defeituosas ($k = 2$).

A probabilidade em uma Distribuição Binomial é calculada pela fórmula:

\begin{equation}
    P(X = k) = \binom{n}{k} \cdot p^k \cdot (1-p)^{n-k}
\end{equation}

Para $n = 5$, $k = 2$, $p = 0{,}2$ e $(1-p) = 0{,}8$:

\begin{equation}
    P(X = 2) = \binom{5}{2} \cdot (0{,}2)^2 \cdot (0{,}8)^{5-2} = 10 \cdot 0{,}04 \cdot 0{,}512 = 0{,}2048
\end{equation}

Como o cálculo de probabilidades da distribuição binomial utiliza operações aritméticas diretas e a combinação simples ($nCr$), certifique-se de retornar ao modo de cálculo geral \textbf{COMP} (\texttt{MODE 1}).

\begin{table}[H]
    \centering
    \begin{tabular}{lll}
        \hline
        \textbf{Etapa} & \textbf{Entrada} & \textbf{Resultado} \\ \hline
        Retornar ao modo COMP & \texttt{MODE 1} & \textit{Modo Padrão} \\
        Cálculo direto de $P(X=2)$ & \texttt{5 nCr 2 * 0.2 \^{} 2 * 0.8 \^{} ( 5 - 2 ) =} & \texttt{0.2048} \\ \hline
    \end{tabular}
\end{table}

A probabilidade de encontrar exatamente 2 peças defeituosas na amostra é de $20{,}48\%$. O uso de parênteses no expoente composto garante a correta ordem das operações.

\subsubsection*{Exemplo 2 — Média e desvio padrão com tabela de frequências}

As notas de uma turma de 20 alunos em uma avaliação de Estatística estão
resumidas na tabela a seguir. Calcule a média e o desvio padrão populacional
utilizando o Modo SD com entrada por frequência.

\begin{table}[H]
    \centering
    \begin{tabular}{ll}
        \hline
        \textbf{Nota} ($x_i$) & \textbf{Frequência} ($f_i$) \\ \hline
        5 & 2 \\
        6 & 4 \\
        7 & 8 \\
        8 & 4 \\
        9 & 2 \\ \hline
    \end{tabular}
\end{table}

\begin{table}[H]
    \centering
    \begin{tabular}{lll}
        \hline
        \textbf{Etapa} & \textbf{Entrada} & \textbf{Visor} \\ \hline
        Ativar Modo SD & \texttt{MODE 2} & SD \\
        Limpar dados & \texttt{SHIFT CLR 1} & 0 \\
        Nota 5, freq.\ 2 & \texttt{5 ; 2 DT} & n = 2 \\
        Nota 6, freq.\ 4 & \texttt{6 ; 4 DT} & n = 6 \\
        Nota 7, freq.\ 8  & \texttt{7 ; 8 DT} & n = 14 \\
        Nota 8, freq.\ 4 & \texttt{8 ; 4 DT} & n = 18 \\
        Nota 9, freq.\ 2 & \texttt{9 ; 2 DT} & n = 20 \\
        Média $\bar{x}$ & \texttt{SHIFT 2 1 =} & 7 \\
        Desvio padrão populacional $s$ & \texttt{SHIFT 2 2 =} & $\approx 1{,}095$ \\ \hline
    \end{tabular}
\end{table}

O resultado $\bar{x} = 7$ pode ser verificado pelo cálculo da média ponderada
pelas frequências:
%
\begin{equation}
    \bar{x} = \frac{5 \cdot 2 + 6 \cdot 4 + 7 \cdot 8 + 8 \cdot 4 + 9 \cdot 2}{20} = \frac{140}{20} = 7
\end{equation}

Como os dados representam toda a turma, utiliza-se o desvio padrão
populacional $s_1$.

\subsubsection*{Exemplo 3 — Coeficiente de variação a partir do Modo SD}

Utilizando os resultados obtidos no Exemplo~1 ($\bar{x} = 7$ e
$s \approx 1{,}414$), calcule o coeficiente de variação. Armazena-se a
média em uma variável para reutilização no cálculo:

\begin{table}[H]
    \centering
    \begin{tabular}{lll}
        \hline
        \textbf{Etapa} & \textbf{Entrada} & \textbf{Resultado} \\ \hline
        Recuperar $\bar{x}$ & \texttt{SHIFT 2 1 =} & 7 \\
        Armazenar em A & \texttt{STO ALPHA A} & 7 \\
        Recuperar $s$ & \texttt{SHIFT 2 3 =} & $\approx 1{,}414$ \\
        Coeficiente de variação & \texttt{Ans / RCL ALPHA A * 100 =}  & $\approx 20{,}2$ \\ \hline
    \end{tabular}
\end{table}

O resultado corresponde à fórmula:
%
\begin{equation}
    CV = \frac{s}{\bar{x}} \times 100 = \frac{1{,}414}{7} \times 100 \approx 20{,}2\%
\end{equation}
